\documentclass[11pt,a4paper]{article}

\usepackage[margin=1in]{geometry}
\usepackage{amsmath,amssymb,amsthm}
\usepackage{booktabs}
\usepackage{multirow}
\usepackage{microtype}
\usepackage{graphicx}
\emergencystretch=4em
\usepackage{hyperref}
\hypersetup{colorlinks=true,linkcolor=black,urlcolor=blue,citecolor=black}

\theoremstyle{plain}
\newtheorem{proposition}{Proposition}
\newtheorem{lemma}[proposition]{Lemma}
\newtheorem{corollary}[proposition]{Corollary}
\theoremstyle{definition}
\newtheorem{remark}[proposition]{Remark}

\newcommand{\R}{\mathbb{R}}
\newcommand{\E}{\mathcal{I}}
\newcommand{\tr}{\operatorname{tr}}
\newcommand{\diag}{\operatorname{diag}}
\newcommand{\spec}{\operatorname{spec}}
\newcommand{\supp}{\operatorname{supp}}
\newcommand{\vecop}{\operatorname{vec}}
\newcommand{\cond}{\kappa}

\title{Closed-Form Observed Information for Coordinate-Based Meta-Regression}
\author{Derivations, equivalence proofs, and the structure of the covariance step}
\date{25 August 2026}

\begin{document}
\maketitle

\begin{abstract}
NiMARE's coordinate-based meta-regression (CBMR) computes regression standard errors by
differentiating the log-likelihood twice with \texttt{torch.func.hessian}. Forward-over-reverse
mode pushes every tangent through the likelihood at once, so the intermediate has shape
$(\text{parameters} \times \text{patterns} \times \text{voxels})$; at $7{,}329$ coefficients,
$21$ patterns and $8{,}444$ voxels that is $9.7$~GB, which is an out-of-memory failure rather
than a slow fit.

We derive the observed information in closed form for all three of NiMARE's observation
distributions, prove each derivation equivalent to the autodiff Hessian, and characterise the
sparsity the design induces in it. A single lemma covers all three: the second differential of
the log-likelihood with respect to the log-intensity is either diagonal in the voxels (Poisson,
Negative Binomial) or diagonal-plus-rank-one (Clustered Negative Binomial), and the chain rule
through the log-bilinear predictor turns each into $n_{\text{patterns}}$ symmetric rank-updates
of width $n_{\text{bases}}$. No parameter-sized intermediate is ever formed. We then separate
two questions that are easily conflated: which \emph{likelihood} admits a closed-form Hessian
(all three), and which \emph{designs} make that Hessian block diagonal (a strictly narrower
class), and we give the generalisations that recover the fast covariance step outside it.
\end{abstract}

\tableofcontents

\section{Setting and notation}

\subsection{The pattern representation}

CBMR models the foci of $n$ experiments over $V$ voxels as a point process whose log-intensity
is bilinear in a spatial spline basis and an experiment-level design. Write

\begin{itemize}
  \item $B \in \R^{V \times K}$ for the spline basis, $b_v^\top$ its $v$-th row;
  \item $L \in \R^{P \times C}$ for the \emph{loadings}: the distinct rows of the experiment-level
        blocks of every spatial term, stacked. $P$ is the number of distinct spatial loadings
        (``patterns''), $C$ the number of spatial columns;
  \item $\pi : \{1,\dots,n\} \to \{1,\dots,P\}$ for the pattern of each experiment, and
        $\E_p = \pi^{-1}(p)$;
  \item $\beta \in \R^{C \times K}$ for the spatial coefficients;
  \item $G \in \R^{n \times Q}$ for the non-spatial (``global'') design block and
        $\gamma \in \R^{Q}$ for its coefficients.
\end{itemize}

The log-intensity of pattern $p$ at voxel $v$, and the scalar moderator effect of experiment
$i$, are
\begin{equation}\label{eq:predictor}
  S_{pv} \;=\; \sum_{c=1}^{C} L_{pc} \sum_{k=1}^{K} \beta_{ck} B_{vk},
  \qquad
  m_i \;=\; \sum_{q=1}^{Q} G_{iq}\gamma_q ,
\end{equation}
so that $S = (L\beta)B^\top$ and $m = G\gamma$. We write $u_{pv} = e^{S_{pv}}$ and
$e_i = e^{m_i}$ throughout. Foci counts are $y_{iv}$, with pattern marginals and experiment
totals
\begin{equation}
  Y_{pv} \;=\; \sum_{i \in \E_p} y_{iv},
  \qquad
  y_{i\cdot} \;=\; \sum_{v=1}^{V} y_{iv}.
\end{equation}

The parameter vector is $\phi = (\vecop(\beta), \gamma) \in \R^{CK + Q}$, with $\beta$ flattened
row-major so that $\beta_{ck}$ occupies index $cK + k$. Overdispersion parameters $\theta_p$,
where the distribution has them, are held fixed at their fitted values --- this is what
\texttt{CBMRModel.information\_matrix} does, and it is what CBMR has always reported regression
standard errors from. Every derivation below is therefore a differentiation in $\phi$ only,
with $\theta$ a constant.

\subsection{The two facts the whole paper rests on}

Both parts of \eqref{eq:predictor} are \emph{linear} in the parameters. Hence
\begin{equation}\label{eq:linear}
  \frac{\partial S_{pv}}{\partial \beta_{ck}} = L_{pc} B_{vk},
  \qquad
  \frac{\partial^2 S_{pv}}{\partial \beta_{ck}\,\partial \beta_{c'k'}} = 0,
  \qquad
  \frac{\partial m_i}{\partial \gamma_q} = G_{iq},
  \qquad
  \frac{\partial^2 m_i}{\partial \gamma_q\,\partial \gamma_{q'}} = 0 .
\end{equation}
The second derivatives vanish identically. Everything that follows is the chain rule applied to
this fact, so the only distribution-specific work is the second differential of the
log-likelihood with respect to $(S, \gamma)$.

\section{A master lemma}

Every CBMR log-likelihood in NiMARE decomposes over patterns,
\begin{equation}\label{eq:decomposition}
  -\log \mathcal{L}(\phi) \;=\; \sum_{p=1}^{P} f_p\bigl(S_{p\cdot},\, \gamma\bigr) ,
\end{equation}
because experiments sharing a pattern share a log-intensity map, and no term couples two
patterns. This is the representation \texttt{predictor.py} is built on.

\begin{lemma}[Chain rule through a log-bilinear predictor]\label{lem:master}
Let $F(\phi) = \sum_p f_p(S_{p\cdot}, \gamma)$ with $S$ as in \eqref{eq:predictor} and each $f_p$
twice differentiable. Write
\[
  \bigl(\Sigma_p\bigr)_{vv'} = \frac{\partial^2 f_p}{\partial S_{pv} \partial S_{pv'}},
  \qquad
  \bigl(\tau_p\bigr)_{vq} = \frac{\partial^2 f_p}{\partial S_{pv} \partial \gamma_q},
  \qquad
  \bigl(\Xi_p\bigr)_{qq'} = \frac{\partial^2 f_p}{\partial \gamma_q \partial \gamma_{q'}} .
\]
Then the Hessian of $F$ has blocks
\begin{align}
  \frac{\partial^2 F}{\partial \beta_{ck}\,\partial \beta_{c'k'}}
    &= \sum_{p=1}^{P} L_{pc} L_{pc'} \, \bigl(B^\top \Sigma_p B\bigr)_{kk'},
      \label{eq:master-bb}\\
  \frac{\partial^2 F}{\partial \beta_{ck}\,\partial \gamma_q}
    &= \sum_{p=1}^{P} L_{pc} \, \bigl(B^\top \tau_p\bigr)_{kq},
      \label{eq:master-bg}\\
  \frac{\partial^2 F}{\partial \gamma_q\,\partial \gamma_{q'}}
    &= \sum_{p=1}^{P} \bigl(\Xi_p\bigr)_{qq'} .
      \label{eq:master-gg}
\end{align}
\end{lemma}

\begin{proof}
By \eqref{eq:linear}, $\partial S_{pv}/\partial\beta_{ck} = L_{pc}B_{vk}$ and the second
derivative of $S$ vanishes, so the chain rule for a composition with an affine map contributes
no curvature term:
\[
  \frac{\partial^2 f_p}{\partial \beta_{ck} \partial \beta_{c'k'}}
  = \sum_{v,v'} \frac{\partial^2 f_p}{\partial S_{pv}\partial S_{pv'}}
      \frac{\partial S_{pv}}{\partial\beta_{ck}} \frac{\partial S_{pv'}}{\partial\beta_{c'k'}}
  = L_{pc}L_{pc'} \sum_{v,v'} (\Sigma_p)_{vv'} B_{vk}B_{v'k'} ,
\]
which is \eqref{eq:master-bb}. The same argument with one factor replaced by
$\partial \gamma_q/\partial\gamma_q = 1$ gives \eqref{eq:master-bg}, and with both replaced,
\eqref{eq:master-gg}. Summing over $p$ completes the proof.
\end{proof}

Lemma~\ref{lem:master} reduces each distribution to a single question: \emph{what is $\Sigma_p$?}
The answer determines the cost. If $\Sigma_p = \diag(w_p)$ then
$B^\top \Sigma_p B = B^\top \diag(w_p) B$ costs $O(VK^2)$ and allocates one $V \times K$
temporary. If $\Sigma_p$ is diagonal plus rank one, one extra outer product of a
$K$-vector is added. In neither case is anything of size $CK$ formed.

\begin{remark}
This is exactly what forward-over-reverse autodiff fails to exploit. \texttt{jacfwd} carries one
tangent per parameter through the evaluation, materialising an array of shape $(CK, P, V)$,
because it cannot know that the map $\beta \mapsto S$ is affine with a rank structure that
collapses the parameter axis.
\end{remark}

\section{Poisson}

\subsection{The likelihood}

Dropping the parameter-free $-\sum \log y!$, NiMARE's Poisson log-likelihood on marginals is
\begin{equation}\label{eq:poisson-ll}
  \log\mathcal{L} \;=\; \sum_{p}\sum_{v} Y_{pv} S_{pv}
  \;+\; \sum_{i} y_{i\cdot}\, m_i
  \;-\; \sum_{p} \Bigl(\sum_v u_{pv}\Bigr) T_p,
  \qquad T_p = \sum_{i \in \E_p} e_i .
\end{equation}

\begin{proposition}[Poisson observed information]\label{prop:poisson}
With $U_{pq} = \sum_{i\in\E_p} e_i G_{iq}$, $a_{pk} = \sum_v u_{pv}B_{vk}$ and
$E_p = \sum_v u_{pv}$, the Hessian of $-\log\mathcal{L}$ is
\begin{align}
  H_{(ck),(c'k')} &= \sum_p L_{pc}L_{pc'} \bigl(B^\top \diag(T_p u_{p\cdot}) B\bigr)_{kk'},\\
  H_{(ck),q} &= \sum_p L_{pc}\, a_{pk} U_{pq},\\
  H_{q,q'} &= \sum_i e_i E_{\pi(i)} G_{iq}G_{iq'} .
\end{align}
In particular $H$ does not depend on the foci.
\end{proposition}

\begin{proof}
Write \eqref{eq:poisson-ll} in the form \eqref{eq:decomposition} with
$f_p = -\sum_v Y_{pv}S_{pv} - \sum_{i\in\E_p} y_{i\cdot}m_i + E_p(S_{p\cdot})\,T_p(\gamma)$.
The first two terms are linear in $\phi$ and so contribute nothing to any second derivative;
this is the whole of the data dependence, which proves the final claim. For the remainder,
$\partial E_p/\partial S_{pv} = u_{pv}$ and $\partial^2 E_p/\partial S_{pv}\partial S_{pv'} =
u_{pv}\delta_{vv'}$, while $\partial T_p/\partial\gamma_q = U_{pq}$ and
$\partial^2 T_p / \partial\gamma_q\partial\gamma_{q'} = \sum_{i\in\E_p} e_i G_{iq}G_{iq'}$.
Hence $\Sigma_p = T_p \diag(u_{p\cdot})$, $(\tau_p)_{vq} = u_{pv}U_{pq}$ and
$(\Xi_p)_{qq'} = E_p \sum_{i \in \E_p} e_iG_{iq}G_{iq'}$. Substituting into
Lemma~\ref{lem:master} and collecting the $\gamma\gamma$ sum over experiments gives the three
displays.
\end{proof}

\begin{remark}
The foci-independence is the log link's canonical-parameter property: for a GLM with canonical
link the observed and expected information coincide. It is a useful test of an implementation,
since a Poisson Hessian that responds to the data is wrong; the test suite checks this.
\end{remark}

\section{Negative Binomial}

\subsection{Reparameterisation}

NiMARE's negative binomial is a moment-matched approximation to the sum, over the experiments of
a pattern, of negative binomial variables. With overdispersion $\theta_p$ and
\begin{equation}
  s_{1} = \sum_{i\in\E_p} e_i, \qquad s_{2} = \sum_{i\in\E_p} e_i^2 ,
\end{equation}
it uses shape $r = s_1^2/(\theta_p s_2)$ and success probability
$\pi_v = \bigl(1 + s_1/(\theta_p u_{pv} s_2)\bigr)^{-1}$. The derivation becomes tractable under
the change of variables
\begin{equation}\label{eq:nb-RA}
  A_p \;=\; \frac{s_1}{\theta_p s_2}, \qquad R_p \;=\; \frac{s_1^2}{\theta_p s_2} \;=\; s_1 A_p ,
\end{equation}
which gives $\pi_v = u_{pv}/(u_{pv}+A_p)$ and $1-\pi_v = A_p/(u_{pv}+A_p)$.

\begin{lemma}\label{lem:nb-psi}
Up to an additive constant depending only on the data,
\begin{equation}\label{eq:psi}
\begin{split}
  f_p \;=\; \Psi(R_p, A_p, S_{p\cdot}) \;=\;{}& -\sum_v \log\Gamma(Y_{pv}+R_p)
     + V\log\Gamma(R_p) - VR_p\log A_p\\
     &{}+ \sum_v (R_p+Y_{pv})\log(u_{pv}+A_p) - \sum_v Y_{pv}S_{pv}.
\end{split}
\end{equation}
\end{lemma}

\begin{proof}
Substituting $\log(1-\pi_v) = \log A_p - \log(u_{pv}+A_p)$ and
$\log \pi_v = S_{pv} - \log(u_{pv}+A_p)$ into NiMARE's summand
$\log\Gamma(Y+r) - \log\Gamma(Y+1) - \log\Gamma(r) + r\log(1-\pi) + Y\log\pi$ and negating gives
\eqref{eq:psi} plus $\sum_v \log\Gamma(Y_{pv}+1)$, which is free of $\phi$.
\end{proof}

The value of \eqref{eq:psi} is that it separates the three ways the parameters enter: the special
functions see only $R_p$, the voxels only $S_{p\cdot}$, and the moderators only $(R_p, A_p)$.

\begin{proposition}[Negative binomial observed information]\label{prop:nb}
Write $d_{pv} = u_{pv}+A_p$. Then
\begin{equation}\label{eq:nb-sigma}
  \Sigma_p = \diag(w_{p\cdot}), \qquad
  w_{pv} \;=\; \frac{(R_p + Y_{pv})\, u_{pv} A_p}{d_{pv}^{2}} ,
\end{equation}
and, with $\partial_R := \partial \Psi/\partial R$ etc.,
\begin{align}
  \frac{\partial^2 \Psi}{\partial S_{pv}\partial R_p} = \frac{u_{pv}}{d_{pv}},
  &\qquad
  \frac{\partial^2 \Psi}{\partial S_{pv}\partial A_p} = -\frac{(R_p+Y_{pv})u_{pv}}{d_{pv}^{2}},
    \label{eq:nb-cross}\\
  \Psi_{RR} = -\sum_v \psi_1(Y_{pv}+R_p) + V\psi_1(R_p),
  &\qquad
  \Psi_{RA} = -\frac{V}{A_p} + \sum_v \frac{1}{d_{pv}},\notag\\
  \Psi_{AA} = \frac{VR_p}{A_p^2} - \sum_v \frac{R_p+Y_{pv}}{d_{pv}^{2}}.&
    \label{eq:nb-second}
\end{align}
where $\psi_1$ is the trigamma function. The $\beta\gamma$ and $\gamma\gamma$ blocks follow by
the chain rule through $R_p(\gamma), A_p(\gamma)$:
\begin{align}
  \frac{\partial^2 f_p}{\partial \beta_{ck}\partial\gamma_q}
   &= L_{pc}\Bigl[\bigl(B^\top \tfrac{u_{p\cdot}}{d_{p\cdot}}\bigr)_k \partial_q R_p
      - \bigl(B^\top \tfrac{(R_p+Y_{p\cdot})u_{p\cdot}}{d_{p\cdot}^2}\bigr)_k \partial_q A_p
      \Bigr], \label{eq:nb-bg}\\
  \frac{\partial^2 f_p}{\partial \gamma_q\partial\gamma_{q'}}
   &= \Psi_{RR}\,\partial_q R_p\, \partial_{q'} R_p
    + \Psi_{RA}\bigl(\partial_q R_p\, \partial_{q'} A_p + \partial_{q'} R_p\, \partial_q A_p\bigr)
    + \Psi_{AA}\, \partial_q A_p\, \partial_{q'} A_p\notag\\
   &\qquad{}+ \Psi_{R}\, \partial^2_{qq'} R_p + \Psi_{A}\, \partial^2_{qq'} A_p.
    \label{eq:nb-gg}
\end{align}
\end{proposition}

\begin{proof}
Differentiate \eqref{eq:psi} in $S_{pv}$, using $\partial u_{pv}/\partial S_{pv} = u_{pv}$:
\[
  \frac{\partial \Psi}{\partial S_{pv}} = \frac{(R_p+Y_{pv})u_{pv}}{d_{pv}} - Y_{pv},
\]
\[
  \frac{\partial^2 \Psi}{\partial S_{pv}\partial S_{pv'}}
    = (R_p+Y_{pv})\,\frac{u_{pv}(u_{pv}+A_p) - u_{pv}^2}{d_{pv}^2}\,\delta_{vv'}
    = \frac{(R_p+Y_{pv})\,u_{pv}A_p}{d_{pv}^2}\,\delta_{vv'},
\]
which is \eqref{eq:nb-sigma}; the Kronecker delta appears because \eqref{eq:psi} is a sum of
per-voxel terms. Differentiating $\partial\Psi/\partial S_{pv}$ in $R_p$ and in $A_p$ gives
\eqref{eq:nb-cross}, and differentiating $\Psi_R, \Psi_A$ gives \eqref{eq:nb-second}. Since
$\gamma$ enters $\Psi$ only through the pair $(R_p, A_p)$, the multivariate chain rule for a
composition $\gamma \mapsto (R,A) \mapsto \Psi$ yields \eqref{eq:nb-bg} and \eqref{eq:nb-gg}, the
last two terms of \eqref{eq:nb-gg} being the curvature of the inner map.
\end{proof}

\begin{lemma}[Derivatives of the reparameterisation]\label{lem:nb-RA-derivs}
Let $U_q = \sum_{i\in\E_p} e_iG_{iq}$, $W_q = \sum_{i\in\E_p} e_i^2G_{iq}$,
$U^{(2)}_{qq'} = \sum_{i\in\E_p} e_iG_{iq}G_{iq'}$,
$W^{(2)}_{qq'} = \sum_{i\in\E_p} e_i^2G_{iq}G_{iq'}$, and
\[
  \ell_q = \frac{U_q}{s_1},\quad
  h_q = \frac{2W_q}{s_2},\quad
  \ell^{(2)}_{qq'} = \frac{U^{(2)}_{qq'}}{s_1} - \frac{U_qU_{q'}}{s_1^2},\quad
  h^{(2)}_{qq'} = \frac{4W^{(2)}_{qq'}}{s_2} - \frac{4W_qW_{q'}}{s_2^2} .
\]
Then, with $\delta^A_q = \ell_q - h_q$ and $\delta^R_q = 2\ell_q - h_q$,
\[
  \partial_q A_p = A_p\,\delta^A_q, \qquad
  \partial_q R_p = R_p\,\delta^R_q,
\]
\[
  \partial^2_{qq'} A_p = A_p\bigl(\delta^A_q\delta^A_{q'} + \ell^{(2)}_{qq'} - h^{(2)}_{qq'}\bigr),
  \qquad
  \partial^2_{qq'} R_p = R_p\bigl(\delta^R_q\delta^R_{q'} + 2\ell^{(2)}_{qq'} - h^{(2)}_{qq'}\bigr).
\]
\end{lemma}

\begin{proof}
$\log A_p = \log s_1 - \log s_2 - \log\theta_p$ and $\log R_p = 2\log s_1 - \log s_2 -
\log\theta_p$. Since $\partial_q s_1 = U_q$ and $\partial_q s_2 = 2W_q$, we get
$\partial_q \log s_1 = \ell_q$ and $\partial_q \log s_2 = h_q$; differentiating again,
$\partial^2_{qq'}\log s_1 = \ell^{(2)}_{qq'}$ and $\partial^2_{qq'} \log s_2 = h^{(2)}_{qq'}$
(the factor $4$ arises because $\partial^2_{qq'} s_2 = 4W^{(2)}_{qq'}$). Finally
$\partial_q X = X \partial_q \log X$ and
$\partial^2_{qq'} X = X\bigl(\partial_q\log X\, \partial_{q'}\log X + \partial^2_{qq'}\log X\bigr)$
for any positive $X$.
\end{proof}

\begin{corollary}[Poisson as a limit]
As $\theta_p \to 0$ with $\gamma$ fixed, $A_p \to \infty$ and $R_p = s_1 A_p \to\infty$, and
\[
  w_{pv} = \frac{(R_p+Y_{pv})u_{pv}A_p}{(u_{pv}+A_p)^2}
   \;\longrightarrow\; s_1 u_{pv} \;=\; T_p u_{pv},
\]
recovering Proposition~\ref{prop:poisson}. In particular the negative binomial weight is
foci-dependent at any $\theta_p > 0$ and the dependence vanishes only in the Poisson limit.
\end{corollary}

\section{Clustered Negative Binomial}

\subsection{The likelihood}

Here the latent effect is a property of the experiment rather than of the voxel, so the spatial
coefficients enter only through the \emph{scalar} $E_p = \sum_v u_{pv}$. With precision
$\nu_p = 1/\theta_p$ and $n_p = |\E_p|$,
\begin{equation}\label{eq:cnb-ll}
\begin{split}
  \log\mathcal{L}_p ={}& n_p \nu_p \log \nu_p - n_p\log\Gamma(\nu_p)
   + \sum_{i\in\E_p}\log\Gamma(y_{i\cdot}+\nu_p)\\
   &{}- \sum_{i\in\E_p}(y_{i\cdot}+\nu_p)\log\bigl(E_p e_i + \nu_p\bigr)
   + \sum_v Y_{pv}S_{pv} + \sum_{i\in\E_p} y_{i\cdot} m_i .
\end{split}
\end{equation}

\begin{proposition}[Clustered negative binomial observed information]\label{prop:cnb}
Let $D_i = E_{\pi(i)} e_i + \nu_{\pi(i)}$, $a_{pk} = \sum_v u_{pv}B_{vk}$ and
\[
  g'_p = \sum_{i\in\E_p} \frac{(y_{i\cdot}+\nu_p)e_i}{D_i},
  \qquad
  g''_p = -\sum_{i\in\E_p} \frac{(y_{i\cdot}+\nu_p)e_i^2}{D_i^2} .
\]
Then $\Sigma_p$ is diagonal plus rank one,
\begin{equation}\label{eq:cnb-sigma}
  \Sigma_p \;=\; g'_p \diag(u_{p\cdot}) \;+\; g''_p\, u_{p\cdot} u_{p\cdot}^\top ,
\end{equation}
and consequently
\begin{align}
  \frac{\partial^2 f_p}{\partial\beta_{ck}\partial\beta_{c'k'}}
   &= L_{pc}L_{pc'}\Bigl[g'_p\bigl(B^\top\diag(u_{p\cdot})B\bigr)_{kk'}
      + g''_p\, a_{pk}a_{pk'}\Bigr],\\
  \frac{\partial^2 f_p}{\partial\beta_{ck}\partial\gamma_q}
   &= L_{pc}\, a_{pk}\, \sum_{i\in\E_p}\frac{(y_{i\cdot}+\nu_p)\nu_p e_i G_{iq}}{D_i^2},\\
  \frac{\partial^2 f_p}{\partial\gamma_q\partial\gamma_{q'}}
   &= \sum_{i\in\E_p}\frac{(y_{i\cdot}+\nu_p)\nu_p E_p e_i}{D_i^2}\,G_{iq}G_{iq'} .
\end{align}
Only the experiment totals $y_{i\cdot}$ appear; the voxelwise marginals $Y_{pv}$ enter
\eqref{eq:cnb-ll} linearly and so vanish.
\end{proposition}

\begin{proof}
Let $\Gamma_p(E,\gamma) = \sum_{i\in\E_p}(y_{i\cdot}+\nu_p)\log(Ee_i+\nu_p)$, so that
$f_p = \Gamma_p(E_p, \gamma)$ up to terms linear in $\phi$ and terms free of $\phi$. Since
$\partial E_p/\partial S_{pv} = u_{pv}$ and
$\partial^2 E_p/\partial S_{pv}\partial S_{pv'} = u_{pv}\delta_{vv'}$,
\[
  \frac{\partial^2 f_p}{\partial S_{pv}\partial S_{pv'}}
   = \Gamma_p''(E_p)\,u_{pv}u_{pv'} + \Gamma_p'(E_p)\,u_{pv}\delta_{vv'} ,
\]
and $\Gamma_p' = g'_p$, $\Gamma_p'' = g''_p$ by direct differentiation, giving
\eqref{eq:cnb-sigma}. Applying Lemma~\ref{lem:master} and using
$B^\top(u u^\top)B = (B^\top u)(B^\top u)^\top = a_pa_p^\top$ gives the $\beta\beta$ block. For
the cross block, $\partial_q \bigl[e_i/D_i\bigr] = \nu_p e_i G_{iq}/D_i^2$, whence
$\partial_q \Gamma'_p = \sum_i (y_{i\cdot}+\nu_p)\nu_pe_iG_{iq}/D_i^2$ and the
$\beta$-derivative contributes the factor $L_{pc}a_{pk}$. For the $\gamma\gamma$ block,
$\partial_q \Gamma_p = \sum_i (y_{i\cdot}+\nu_p)E_pe_iG_{iq}/D_i$ and
$\partial_{q'}\bigl[E_pe_i/D_i\bigr] = E_pe_i\nu_pG_{iq'}/D_i^2$.
\end{proof}

\begin{remark}
Proposition~\ref{prop:cnb} is the only one of the three whose $\Sigma_p$ is not diagonal, and it
is nonetheless the cheapest to prove: with $\nu_p$ fixed, every special function in
\eqref{eq:cnb-ll} depends on constants alone.
\end{remark}

\section{Equivalence}

Two independent checks establish that the propositions describe the same matrix NiMARE computes.

\subsection{Symbolic}

\texttt{scripts/symbolic\_nb.py} and \texttt{scripts/symbolic\_cnb.py} construct the likelihood
exactly as \texttt{nimare/meta/cbmr/distributions.py} writes it, differentiate it with
\texttt{sympy} in the primitive coefficients, subtract the claimed closed form, and simplify.
Every residual is the symbol $0$; nothing numerical is involved, so these are proofs on a
symbolic instance rather than evidence at a sample point.

For the clustered case the whole composite Hessian closes in one pass. For the negative binomial
the composite is intractable for \texttt{simplify}, so the proof is decomposed exactly as the
argument above is: (i) NiMARE's summand equals $\Psi(R,A,S)$ of Lemma~\ref{lem:nb-psi}; (ii) the
second partials of $\Psi$ are \eqref{eq:nb-sigma}--\eqref{eq:nb-second}; (iii) the derivatives of
$R,A$ are those of Lemma~\ref{lem:nb-RA-derivs}. Composition is then
Lemma~\ref{lem:master} together with the chain rule, which is a theorem.

\begin{remark}[A subtlety worth recording]
A first attempt verified (i) by eliminating $s_1$ from NiMARE's expression via $s_1 = A\theta s_2$
and comparing with $\Psi$. This reports a spurious mismatch whose residual contains
$A^2 s_2 \theta$. That quantity \emph{is} $R$: given $(s_1,s_2,\theta)$, $R$ and $A$ are not
independent, and eliminating only one of the two relations leaves the expressions in different
but equal forms. Substituting both definitions into $\Psi$ instead closes the identity.
\end{remark}

\subsection{Numerical}

\texttt{scripts/test\_cbmr\_analytic.py} and \texttt{scripts/test\_cbmr\_designs.py} compare
the implementations --- vectorisation, pattern assembly, loadings scatter --- against
\texttt{CBMRModel.information\_matrix}, i.e. against \texttt{torch.func.hessian} itself.
Table~\ref{tab:designs} reports the design coverage. Two further tests guard the structural
claims: the Poisson Hessian must be unchanged when the foci are doubled
(Proposition~\ref{prop:poisson}), and the other two must move, so a foci-independent
implementation of either cannot pass silently.

\section{Structure of the information matrix}

\subsection{Sparsity from the design}

\begin{proposition}[Column separation]\label{prop:blocks}
Let $\supp(p) = \{c : L_{pc} \neq 0\}$ and let $\sim$ be the transitive closure of the relation
$c \sim c' \iff \exists p:\ c,c' \in \supp(p)$. If $Q = 0$ and $c \not\sim c'$ then
$H_{(ck),(c'k')} = 0$ for all $k,k'$. Consequently $H$ is block diagonal over the connected
components of $\sim$, each block of size $|{\cdot}| \cdot K$.
\end{proposition}

\begin{proof}
Immediate from \eqref{eq:master-bb}: each summand carries the factor $L_{pc}L_{pc'}$, which
vanishes unless both $c$ and $c'$ lie in $\supp(p)$; if they do, then $c \sim c'$.
\end{proof}

Proposition~\ref{prop:blocks} holds for \emph{any} $\Sigma_p$ and hence for all three
distributions --- and indeed for any likelihood of the form \eqref{eq:decomposition}. It is a
statement about the design, not about the observation model. Note the entries are
\emph{structurally} zero, not small: no rounding decision is involved.

\subsection{Which designs separate}

The hypothesis $Q=0$ is not cosmetic, and neither is the connectivity condition. A cell-means
spatial factor \texttt{s(a)} assigns each pattern exactly one column, so $\sim$ is the identity
and there are $C$ blocks; so does a full interaction \texttt{s(a:b)}. Every other construction
the CBMR grammar admits couples the columns:
\begin{itemize}
  \item a sum-to-zero term \texttt{sz(a)} spans all levels in every row, so $\supp(p)$ is the
        whole term and one component absorbs everything;
  \item a spatial covariate \texttt{s(x)} loads alongside whatever factor accompanies it;
  \item any global moderator makes $Q > 0$, and the cross block \eqref{eq:master-bg} is generally
        dense, coupling every spatial column through $\gamma$.
\end{itemize}
Table~\ref{tab:designs} makes this concrete. The practical consequence is that the closed-form
Hessian is general while block-diagonal inversion is not, and the two must be gated separately.

\subsection{Bordered structure}\label{sec:bordered}

When the spatial columns separate but $Q>0$, $H$ is not block diagonal but is \emph{bordered}
block diagonal,
\begin{equation}\label{eq:bordered}
  H = \begin{pmatrix} D & C \\ C^\top & E\end{pmatrix},
  \qquad D = \operatorname{blkdiag}(D_1,\dots,D_M) \in \R^{CK\times CK},\quad C \in \R^{CK\times Q},
\end{equation}
with $Q$ typically between $1$ and $5$.

\begin{proposition}[Bordered inverse]\label{prop:schur}
If $D$ and the Schur complement $\Sigma = E - C^\top D^{-1}C$ are nonsingular, then
\begin{equation}
  H^{-1} =
  \begin{pmatrix}
    D^{-1} + D^{-1}C\,\Sigma^{-1}C^\top D^{-1} & -D^{-1}C\,\Sigma^{-1}\\[2pt]
    -\Sigma^{-1}C^\top D^{-1} & \Sigma^{-1}
  \end{pmatrix},
\end{equation}
computable in $O\bigl(\sum_j |D_j|^3\bigr) + O(CKQ\cdot\max_j|D_j|) + O(Q^3)$ plus the
$O((CK)^2Q)$ cost of forming the outer product, rather than $O((CK+Q)^3)$.
\end{proposition}

\begin{proof}
Standard block inversion; verify by multiplying out $H H^{-1} = I$ blockwise, using
$E - C^\top D^{-1} C = \Sigma$.
\end{proof}

Proposition~\ref{prop:schur} extends the fast \emph{inverse} to designs of the form
``cell-means factor plus global moderators''. It does not extend the fast \emph{condition
number}: the spectrum of \eqref{eq:bordered} is not the union of the spectra of $D$ and
$\Sigma$, and no interlacing bound recovers the value.

\section{Condition number and inversion}

\texttt{CBMRModel.covariance} computes $\cond_2(H) = \sigma_{\max}/\sigma_{\min}$ with
\texttt{np.linalg.cond}, which forms a full singular value decomposition, and then inverts with
\texttt{np.linalg.inv}. At the $21$-group shape these cost $163$~s and $14$~s respectively, so
the diagnostic dominates the quantity it is diagnosing. Three generalisations are available, in
increasing order of applicability.

\begin{proposition}[Block-diagonal spectrum]\label{prop:spectrum}
If $H = \operatorname{blkdiag}(H_1,\dots,H_M)$ then
$\spec(H) = \bigsqcup_j \spec(H_j)$ and the
singular values are likewise the multiset union. Hence
$\cond_2(H) = \max_j\sigma_{\max}(H_j) \big/ \min_j \sigma_{\min}(H_j)$, exactly.
\end{proposition}

\begin{proof}
$H - \lambda I$ is block diagonal with blocks $H_j - \lambda I$, and the determinant of a block
diagonal matrix is the product of the blocks' determinants, so the characteristic polynomial
factorises. The same argument applied to $H^\top H$ gives the singular values.
\end{proof}

\begin{proposition}[Symmetry suffices]\label{prop:eigh}
$H$ is symmetric by construction. For symmetric $H$ with eigenvalues $\lambda_i$, the singular
values are $|\lambda_i|$, so $\cond_2(H) = \max_i|\lambda_i| / \min_i|\lambda_i|$ and a symmetric
eigensolver may replace the SVD. This requires no hypothesis on the design and no positive
definiteness.
\end{proposition}

\begin{proof}
Write $H = V\Lambda V^\top$ with $V$ orthogonal. Then $H^\top H = V\Lambda^2V^\top$, so the
singular values --- the nonnegative square roots of the eigenvalues of $H^\top H$ --- are
$|\lambda_i|$.
\end{proof}

\begin{proposition}[One factorisation for both]\label{prop:chol}
If $H \succ 0$, a single Cholesky factorisation $H = LL^\top$ serves both purposes: LAPACK's
\texttt{pocon} returns an estimate of $\cond_1(H)$ in $O(n^2)$ from $L$, and \texttt{potri}
returns $H^{-1}$ in $O(n^3/3)$ from the same $L$. Moreover $\cond_1$ and $\cond_2$ are equivalent
within a factor of $n$:
\begin{equation}\label{eq:norm-equiv}
  n^{-1}\cond_1(H) \;\le\; \cond_2(H) \;\le\; n\,\cond_1(H).
\end{equation}
\end{proposition}

\begin{proof}
The factorisation and the two LAPACK routines are standard. For \eqref{eq:norm-equiv}, from
$n^{-1/2}\|M\|_1 \le \|M\|_2 \le n^{1/2}\|M\|_1$ applied to $M = H$ and $M = H^{-1}$,
$\cond_2 = \|H\|_2\|H^{-1}\|_2 \le n^{1/2}\|H\|_1 \cdot n^{1/2}\|H^{-1}\|_1 = n\cond_1$, and
symmetrically.
\end{proof}

\begin{remark}[The bound is not the behaviour]\label{rem:pocon}
Proposition~\ref{prop:chol} is the fastest route by a wide margin, and \eqref{eq:norm-equiv}
makes it sound harmless: a factor of $n$ is four of the sixteen decimal digits. Measured, it is
not harmless. On the three structures of Table~\ref{tab:strategies} the \texttt{pocon} estimate
reports $90.09$ against a true $\cond_2$ of $5.02$, $2.42\times10^{5}$ against $1657.7$, and
$1312.4$ against $4.98$ --- discrepancies of $18\times$, $146\times$ and $264\times$. These are
comfortably inside the bound, which is the point: the bound permits behaviour bad enough to
make the estimate useless as a drop-in for a number the user sees. A fit at
$\cond_2 = 10^{14}$ --- poorly conditioned but still invertible --- would be reported at
$\sim\!10^{16}$ and trip a warning saying its standard errors cannot be trusted.

The conclusion is to split the two uses of the factorisation. The Cholesky is adopted for the
\emph{inverse}, where it agrees with \texttt{np.linalg.inv} to $1.8\times10^{-15}$ and is
$2$--$3\times$ faster. It is rejected for the \emph{condition number}, which comes instead from
Proposition~\ref{prop:eigh}: exact, fully general, and still $3$--$4\times$ faster than the SVD
the stock path forms.
\end{remark}

\subsection{The ladder}\label{sec:ladder}

Collecting Propositions~\ref{prop:schur}--\ref{prop:chol} into a dispatch, most specific rung
first:

\begin{enumerate}
  \item \textbf{blockwise} --- the spatial columns separate and $Q = 0$. Condition number by
        Proposition~\ref{prop:spectrum}, inverse blockwise. Reached only by a lone cell-means
        factor or a lone full interaction.
  \item \textbf{bordered} --- the spatial columns separate but $Q > 0$. Inverse by
        Proposition~\ref{prop:schur}; the condition number cannot follow, so it falls back to
        Proposition~\ref{prop:eigh} on the whole matrix.
  \item \textbf{Cholesky} --- no exploitable structure, but $H \succ 0$. Condition number by
        Proposition~\ref{prop:eigh}, inverse by Cholesky.
  \item \textbf{dense} --- $H \not\succ 0$, so the Cholesky declines. Falls through to
        \texttt{np.linalg.inv}, which either succeeds or raises with the diagnostic the stock
        path would have raised.
\end{enumerate}

Every rung reports the \emph{same} condition number, to the last printed digit; none of them is
an estimate. What varies is only how much work is done to obtain it. Structural claims are
checked against the matrix in hand rather than assumed --- a design that ought to separate but
whose information matrix has an entry outside the implied blocks is inverted densely, with a
warning, rather than silently truncated.

Table~\ref{tab:designs} records which rung each design reaches. The fourth is reached by exactly
one design in the suite, \texttt{\textasciitilde\ sz(a) + sz(b) + sz(a:b)}, which is rank
deficient at that sample size ($\cond_2 \approx 4\times10^{20}$); the Cholesky declining is the
guard working.

\section{Empirical results}

\subsection{Cost of the closed forms}

\begin{table}[htbp]
\centering
\caption{Whole \texttt{covariance()} call at the 21-group diagnosis shape: $7{,}329$
coefficients, $21$ patterns, $8{,}444$ voxels, $349$ bases. ``chunked jacrev'' is
$\texttt{jacrev}(\texttt{jacrev}(f), \text{chunk}=192)$ followed by \texttt{np.linalg.cond}
and \texttt{np.linalg.inv}; stock \texttt{torch.func.hessian} does not run at this size at all,
requiring a $9.7$~GB intermediate. Memory is the increment over the pre-call peak. The last
column is the largest relative disagreement between the two covariance matrices.}
\label{tab:scale}
\begin{tabular}{lrrrrrr}
\toprule
& \multicolumn{2}{c}{chunked jacrev} & \multicolumn{2}{c}{closed form} & & \\
\cmidrule(lr){2-3}\cmidrule(lr){4-5}
distribution & time & memory & time & memory & speed-up & disagreement \\
\midrule
Poisson                    & $205.4$~s & $+2.29$~GB & $1.5$~s & $+0.13$~GB & $139\times$ & $5.3\times10^{-15}$ \\
NegativeBinomial           & $343.8$~s & $+2.74$~GB & $1.5$~s & $+0.00$~GB & $226\times$ & $6.3\times10^{-15}$ \\
ClusteredNegativeBinomial  & $287.3$~s & ---        & $2.3$~s & ---        & $124\times$ & $3.7\times10^{-15}$ \\
\bottomrule
\end{tabular}
\end{table}


\subsection{Design coverage}

\begin{table}[htbp]
\centering
\caption{Design coverage. ``patterns'' is $P$ and ``$\max|\supp(p)|$'' the largest number of
spatial columns any one pattern loads on --- greater than one exactly when the loadings are not
one-hot. ``blocks'' is the number of components of Proposition~\ref{prop:blocks}; ``rung'' is
which strategy of Section~\ref{sec:ladder} that buys. ``disagreement'' is the largest relative
difference between the closed-form Hessian and \texttt{torch.func.hessian}, over the
distributions that accept the design. Every row was checked for all three distributions except
where noted.}
\label{tab:designs}
\resizebox{\textwidth}{!}{%
\begin{tabular}{llrrrlr}
\toprule
design & formula & $P$ & $\max|\supp(p)|$ & blocks & rung & disagreement \\
\midrule
cell-means factor & \texttt{\textasciitilde\ s(a)} & $4$ & $1$ & $4$ & blockwise & $2.7\times10^{-15}$\\
full interaction & \texttt{\textasciitilde\ s(a:b)} & $8$ & $1$ & $8$ & blockwise & $2.7\times10^{-15}$\\
additive main effects & \texttt{\textasciitilde\ sz(a) + sz(b)} & $8$ & $5$ & $1$ & Cholesky & $5.0\times10^{-16}$\\
main effects $+$ interaction & \texttt{\textasciitilde\ sz(a) + sz(b) + sz(a:b)} & $8$ & $12$ & $1$ & dense$^{\dagger}$ & $4.6\times10^{-16}$\\
one global moderator & \texttt{\textasciitilde\ s(a) + x} & $4$ & $1$ & $1$ & bordered & $2.2\times10^{-14}$\\
global moderator interaction & \texttt{\textasciitilde\ s(a) + x * z} & $4$ & $1$ & $1$ & bordered & $9.2\times10^{-14}$\\
voxelwise moderator, continuous & \texttt{\textasciitilde\ s(a) + s(x)} & $96$ & $2$ & $1$ & Cholesky & $2.3\times10^{-16\,\ddagger}$\\
voxelwise moderator, coarse & \texttt{\textasciitilde\ s(a) + s(w)} & $12$ & $2$ & $1$ & Cholesky & $4.1\times10^{-16}$\\
everything at once & \texttt{\textasciitilde\ sz(a) + sz(b) + s(w) + x * z} & $24$ & $6$ & $1$ & Cholesky & $1.2\times10^{-14}$\\
\bottomrule
\end{tabular}}

\vspace{2pt}
\raggedright\footnotesize
$^{\dagger}$ Rank deficient at this sample size (condition number $4\times10^{20}$), so the
Cholesky declines and the ladder falls through to the general inverse. The guard working, not a
gap.\\
$^{\ddagger}$ Poisson only: NiMARE's \texttt{check\_design} refuses the overdispersed
distributions when a continuous spatial covariate leaves every pattern with one experiment. The
same term with a coarse covariate (previous row) is accepted by all three.
\end{table}


\subsection{Condition number and inverse}

\begin{table}[htbp]
\centering
\caption{Condition number and inverse, on synthetic matrices at the 21-group shape with the
three structures real designs produce. ``dense'' is the stock path
(\texttt{np.linalg.cond} $+$ \texttt{np.linalg.inv}). All inverses agree with the dense one to
$\le 4.6\times10^{-15}$ relative. Rows marked $\star$ are what the ladder adopts. The
blockwise row was measured with a per-block SVD; the implementation uses a per-block symmetric
eigensolver, which returns the same values by Propositions~\ref{prop:spectrum}
and~\ref{prop:eigh}.}
\label{tab:strategies}
\resizebox{\textwidth}{!}{%
\begin{tabular}{llrrrr}
\toprule
structure & strategy & $\cond$ reported & $\cond$ time & inverse time & total \\
\midrule
\multirow{4}{*}{\shortstack[l]{block diagonal\\ \texttt{\textasciitilde\ s(a)}, $7329^2$}}
 & dense (SVD $+$ LU)          & $5.0243$              & $157.2$~s & $14.2$~s & $171.4$~s\\
 & \texttt{eigvalsh} $+$ Cholesky & $5.0243$           & $36.5$~s  & $7.6$~s  & $44.1$~s\\
 & Cholesky $+$ \texttt{pocon} & $90.093$              & $2.9$~s   & $5.5$~s  & $8.4$~s\\
 & $\star$ blockwise           & $5.0243$              & $0.4$~s   & $0.2$~s  & $\mathbf{0.7}$~s\\
\midrule
\multirow{4}{*}{\shortstack[l]{bordered\\ \texttt{\textasciitilde\ s(a) + x * z}, $7331^2$}}
 & dense (SVD $+$ LU)          & $1657.7$              & $169.0$~s & $20.7$~s & $189.6$~s\\
 & $\star$ \texttt{eigvalsh} $+$ Schur & $1657.7$      & $48.5$~s  & $4.1$~s  & $\mathbf{52.7}$~s\\
 & Cholesky $+$ \texttt{pocon} & $2.42\times10^{5}$    & $4.7$~s   & $8.7$~s  & $13.3$~s\\
 & blockwise                   & \multicolumn{4}{c}{\emph{not applicable --- the matrix does not separate}}\\
\midrule
\multirow{3}{*}{\shortstack[l]{fully coupled\\ \texttt{\textasciitilde\ sz(a) + sz(b)}, $7329^2$}}
 & dense (SVD $+$ LU)          & $4.9796$              & $161.8$~s & $19.8$~s & $181.7$~s\\
 & $\star$ \texttt{eigvalsh} $+$ Cholesky & $4.9796$   & $52.1$~s  & $15.8$~s & $\mathbf{67.8}$~s\\
 & Cholesky $+$ \texttt{pocon} & $1312.4$              & $4.4$~s   & $8.6$~s  & $13.0$~s\\
\bottomrule
\end{tabular}}
\end{table}


\section{Scope and limitations}

\begin{enumerate}
  \item The derivations hold with the overdispersion parameters $\theta_p$ held fixed. This
        matches what \texttt{information\_matrix} does and what CBMR reports, but it means the
        reported standard errors do not propagate uncertainty in $\theta$. Extending to the joint
        $(\phi,\theta)$ information would require the $\theta$ blocks, which are not derived here.
  \item NiMARE's \texttt{check\_design} refuses the overdispersed distributions when any pattern
        has a single experiment, which excludes a continuous spatial covariate. That is a
        pre-existing restriction of the moment-matching approximation, not of these derivations;
        the same term with a coarse covariate is accepted and handled.
  \item The negative binomial $\gamma\gamma$ block \eqref{eq:nb-gg} is the least numerically
        comfortable part of the implementation: $\Psi_R\,\partial^2R$ and $\Psi_A\,\partial^2A$ are
        large and opposite in sign, and agreement with autodiff degrades from $\sim\!10^{-16}$ to
        $\sim\!10^{-13}$ relative when moderators are present. This is far below any tolerance
        that matters for standard errors but is worth knowing.
  \item The covariance is returned as a dense $(CK+Q)^2$ array because
        \texttt{contrasts.py} indexes it densely. At the $21$-group shape that is $430$~MB, and it
        is the remaining memory floor; exploiting the block structure past the covariance itself
        would require changes inside NiMARE.
\end{enumerate}

\clearpage
% Collects the machine-generated proof fragments into one appendix.
\section{Generated proof fragments}
\label{sec:generated}

Everything below is written by \texttt{run\_proofs.py} from the same sympy objects the proofs
check. Nothing here is transcribed by hand, so the document cannot disagree with the
verification that backs it.

% Generated by cbmr-proofs. Do not edit.
% Every expression below is rendered from the object sympy verified, not transcribed.

\subsection{Poisson}

The log-likelihood on marginals, dropping the parameter-free
$-\sum\log y!$, is $\log\mathcal{L}=\sum_{v}Y_{v}S_{v}+\sum_i y_i m_i-(\sum_v e^{S_v})T$ with
$T=\sum_i e^{m_i}$. The first two terms are linear in the coefficients, so they differentiate
away entirely; that is the whole of the data dependence.


\noindent The expressions the claims are formed from:

\begin{equation*}
  -\log\mathcal{L} = - Y_{0} \left(B_{00} \beta_{0} + B_{01} \beta_{1}\right) - Y_{1} \left(B_{10} \beta_{0} + B_{11} \beta_{1}\right) - Y_{2} \left(B_{20} \beta_{0} + B_{21} \beta_{1}\right) - y_{0} \left(G_{00} \gamma_{0} + G_{01} \gamma_{1}\right) - y_{1} \left(G_{10} \gamma_{0} + G_{11} \gamma_{1}\right) + \left(e^{G_{00} \gamma_{0} + G_{01} \gamma_{1}} + e^{G_{10} \gamma_{0} + G_{11} \gamma_{1}}\right) \left(e^{B_{00} \beta_{0} + B_{01} \beta_{1}} + e^{B_{10} \beta_{0} + B_{11} \beta_{1}} + e^{B_{20} \beta_{0} + B_{21} \beta_{1}}\right)
\end{equation*}

\noindent Residuals, each asserted to vanish:

\begin{center}\small
\begin{tabular}{lll}
\toprule
claim & residual & result\\
\midrule
d2/dbeta0 dbeta0 & $\partial^2_{\beta_0\beta_0} - T\,(B^\top\mathrm{diag}(u)B)_{00}$ & $0$\\
d2/dbeta0 dbeta1 & $\partial^2_{\beta_0\beta_1} - T\,(B^\top\mathrm{diag}(u)B)_{01}$ & $0$\\
d2/dbeta1 dbeta0 & $\partial^2_{\beta_1\beta_0} - T\,(B^\top\mathrm{diag}(u)B)_{10}$ & $0$\\
d2/dbeta1 dbeta1 & $\partial^2_{\beta_1\beta_1} - T\,(B^\top\mathrm{diag}(u)B)_{11}$ & $0$\\
d2/dbeta0 dgamma0 & $\partial^2_{\beta_0\gamma_0} - a_0U_0$ & $0$\\
d2/dbeta0 dgamma1 & $\partial^2_{\beta_0\gamma_1} - a_0U_1$ & $0$\\
d2/dbeta1 dgamma0 & $\partial^2_{\beta_1\gamma_0} - a_1U_0$ & $0$\\
d2/dbeta1 dgamma1 & $\partial^2_{\beta_1\gamma_1} - a_1U_1$ & $0$\\
d2/dgamma0 dgamma0 & $\partial^2_{\gamma_0\gamma_0} - E\sum_i e_iG_{i0}G_{i0}$ & $0$\\
d2/dgamma0 dgamma1 & $\partial^2_{\gamma_0\gamma_1} - E\sum_i e_iG_{i0}G_{i1}$ & $0$\\
d2/dgamma1 dgamma0 & $\partial^2_{\gamma_1\gamma_0} - E\sum_i e_iG_{i1}G_{i0}$ & $0$\\
d2/dgamma1 dgamma1 & $\partial^2_{\gamma_1\gamma_1} - E\sum_i e_iG_{i1}G_{i1}$ & $0$\\
Hessian contains no foci symbol & $\{Y_v, y_i\} \cap \mathrm{free}(H)$ & $0$\\
\bottomrule
\end{tabular}
\end{center}

\noindent\textbf{Verdict:} every residual is exactly zero.

% Generated by cbmr-proofs. Do not edit.
% Every expression below is rendered from the object sympy verified, not transcribed.

\subsection{Negative Binomial}

NiMARE's negative binomial uses shape $r=s_1^2/(\theta s_2)$ and success
probability $\pi_v=(1+s_1/(\theta u_v s_2))^{-1}$. Under the reparameterisation
$A=s_1/(\theta s_2)$, $R=s_1A$ we have $\pi_v=u_v/(u_v+A)$, and the special functions depend on
$R$ alone, the voxels on $S$ alone, and the moderators on $(R,A)$ alone.


\noindent The expressions the claims are formed from:

\begin{equation*}
  \Psi(R, A, S) = - 3 R \log{\left(A \right)} - S_{0} y_{0} - S_{1} y_{1} - S_{2} y_{2} + \left(R + y_{0}\right) \log{\left(A + e^{S_{0}} \right)} + \left(R + y_{1}\right) \log{\left(A + e^{S_{1}} \right)} + \left(R + y_{2}\right) \log{\left(A + e^{S_{2}} \right)} + 3 \operatorname{loggamma}{\left(R \right)} - \operatorname{loggamma}{\left(R + y_{0} \right)} - \operatorname{loggamma}{\left(R + y_{1} \right)} - \operatorname{loggamma}{\left(R + y_{2} \right)}
\end{equation*}

\noindent Residuals, each asserted to vanish:

\begin{center}\small
\begin{tabular}{lll}
\toprule
claim & residual & result\\
\midrule
Psi(R,A,S) equals NiMARE's summand & $\Psi\big|_{R,A} - \big(-\log\mathcal{L} - \textrm{const}\big)$ & $0$\\
d2Psi/dS0\^{}2 & $\partial^2_{S_0}\Psi - \frac{(R+Y_0)u_0A}{(u_0+A)^2}$ & $0$\\
d2Psi/dS0 dR & $\partial^2_{S_0R}\Psi - \frac{u_0}{u_0+A}$ & $0$\\
d2Psi/dS0 dA & $\partial^2_{S_0A}\Psi + \frac{(R+Y_0)u_0}{(u_0+A)^2}$ & $0$\\
d2Psi/dS1\^{}2 & $\partial^2_{S_1}\Psi - \frac{(R+Y_1)u_1A}{(u_1+A)^2}$ & $0$\\
d2Psi/dS1 dR & $\partial^2_{S_1R}\Psi - \frac{u_1}{u_1+A}$ & $0$\\
d2Psi/dS1 dA & $\partial^2_{S_1A}\Psi + \frac{(R+Y_1)u_1}{(u_1+A)^2}$ & $0$\\
d2Psi/dS2\^{}2 & $\partial^2_{S_2}\Psi - \frac{(R+Y_2)u_2A}{(u_2+A)^2}$ & $0$\\
d2Psi/dS2 dR & $\partial^2_{S_2R}\Psi - \frac{u_2}{u_2+A}$ & $0$\\
d2Psi/dS2 dA & $\partial^2_{S_2A}\Psi + \frac{(R+Y_2)u_2}{(u_2+A)^2}$ & $0$\\
Psi\_R & $\Psi_R\ \textrm{as implemented}$ & $0$\\
Psi\_A & $\Psi_A\ \textrm{as implemented}$ & $0$\\
Psi\_RR & $\Psi_{RR}\ \textrm{as implemented}$ & $0$\\
Psi\_RA & $\Psi_{RA}\ \textrm{as implemented}$ & $0$\\
Psi\_AA & $\Psi_{AA}\ \textrm{as implemented}$ & $0$\\
dR/dgamma0 & $\partial_{\gamma_0}R - R(2\ell_0-h_0)$ & $0$\\
dA/dgamma0 & $\partial_{\gamma_0}A - A(\ell_0-h_0)$ & $0$\\
d2R/dgamma0 dgamma0 & $\partial^2_{\gamma_0\gamma_0}R\ \textrm{as implemented}$ & $0$\\
d2A/dgamma0 dgamma0 & $\partial^2_{\gamma_0\gamma_0}A\ \textrm{as implemented}$ & $0$\\
d2R/dgamma0 dgamma1 & $\partial^2_{\gamma_0\gamma_1}R\ \textrm{as implemented}$ & $0$\\
d2A/dgamma0 dgamma1 & $\partial^2_{\gamma_0\gamma_1}A\ \textrm{as implemented}$ & $0$\\
dR/dgamma1 & $\partial_{\gamma_1}R - R(2\ell_1-h_1)$ & $0$\\
dA/dgamma1 & $\partial_{\gamma_1}A - A(\ell_1-h_1)$ & $0$\\
d2R/dgamma1 dgamma0 & $\partial^2_{\gamma_1\gamma_0}R\ \textrm{as implemented}$ & $0$\\
d2A/dgamma1 dgamma0 & $\partial^2_{\gamma_1\gamma_0}A\ \textrm{as implemented}$ & $0$\\
d2R/dgamma1 dgamma1 & $\partial^2_{\gamma_1\gamma_1}R\ \textrm{as implemented}$ & $0$\\
d2A/dgamma1 dgamma1 & $\partial^2_{\gamma_1\gamma_1}A\ \textrm{as implemented}$ & $0$\\
\bottomrule
\end{tabular}
\end{center}

\noindent\textbf{Verdict:} every residual is exactly zero.

% Generated by cbmr-proofs. Do not edit.
% Every expression below is rendered from the object sympy verified, not transcribed.

\subsection{Clustered Negative Binomial}

The latent effect is a property of the experiment rather than of the voxel, so
the spatial coefficients reach the likelihood only through the scalar $E=\sum_v u_v$. Writing
$\Gamma(E)=\sum_i (y_i+\nu)\log(Ee_i+\nu)$, the second differential in $S$ is
$\Gamma''(E)u_vu_{v'}+\Gamma'(E)u_v\delta_{vv'}$ --- diagonal plus rank one.


\noindent The expressions the claims are formed from:

\begin{equation*}
  -\log\mathcal{L}_p = - Y_{0} \left(B_{00} \beta_{0} + B_{01} \beta_{1}\right) - Y_{1} \left(B_{10} \beta_{0} + B_{11} \beta_{1}\right) - Y_{2} \left(B_{20} \beta_{0} + B_{21} \beta_{1}\right) - 2 \nu \log{\left(\nu \right)} - y_{0} \left(G_{00} \gamma_{0} + G_{01} \gamma_{1}\right) - y_{1} \left(G_{10} \gamma_{0} + G_{11} \gamma_{1}\right) + \left(\nu + y_{0}\right) \log{\left(\nu + \left(e^{B_{00} \beta_{0} + B_{01} \beta_{1}} + e^{B_{10} \beta_{0} + B_{11} \beta_{1}} + e^{B_{20} \beta_{0} + B_{21} \beta_{1}}\right) e^{G_{00} \gamma_{0} + G_{01} \gamma_{1}} \right)} + \left(\nu + y_{1}\right) \log{\left(\nu + \left(e^{B_{00} \beta_{0} + B_{01} \beta_{1}} + e^{B_{10} \beta_{0} + B_{11} \beta_{1}} + e^{B_{20} \beta_{0} + B_{21} \beta_{1}}\right) e^{G_{10} \gamma_{0} + G_{11} \gamma_{1}} \right)} + 2 \operatorname{loggamma}{\left(\nu \right)} - \operatorname{loggamma}{\left(\nu + y_{0} \right)} - \operatorname{loggamma}{\left(\nu + y_{1} \right)}
\end{equation*}
\begin{equation*}
  g' = \frac{\left(\nu + y_{0}\right) e^{G_{00} \gamma_{0} + G_{01} \gamma_{1}}}{\nu + \left(e^{B_{00} \beta_{0} + B_{01} \beta_{1}} + e^{B_{10} \beta_{0} + B_{11} \beta_{1}} + e^{B_{20} \beta_{0} + B_{21} \beta_{1}}\right) e^{G_{00} \gamma_{0} + G_{01} \gamma_{1}}} + \frac{\left(\nu + y_{1}\right) e^{G_{10} \gamma_{0} + G_{11} \gamma_{1}}}{\nu + \left(e^{B_{00} \beta_{0} + B_{01} \beta_{1}} + e^{B_{10} \beta_{0} + B_{11} \beta_{1}} + e^{B_{20} \beta_{0} + B_{21} \beta_{1}}\right) e^{G_{10} \gamma_{0} + G_{11} \gamma_{1}}}
\end{equation*}
\begin{equation*}
  g'' = - \frac{\left(\nu + y_{0}\right) e^{2 G_{00} \gamma_{0} + 2 G_{01} \gamma_{1}}}{\left(\nu + \left(e^{B_{00} \beta_{0} + B_{01} \beta_{1}} + e^{B_{10} \beta_{0} + B_{11} \beta_{1}} + e^{B_{20} \beta_{0} + B_{21} \beta_{1}}\right) e^{G_{00} \gamma_{0} + G_{01} \gamma_{1}}\right)^{2}} - \frac{\left(\nu + y_{1}\right) e^{2 G_{10} \gamma_{0} + 2 G_{11} \gamma_{1}}}{\left(\nu + \left(e^{B_{00} \beta_{0} + B_{01} \beta_{1}} + e^{B_{10} \beta_{0} + B_{11} \beta_{1}} + e^{B_{20} \beta_{0} + B_{21} \beta_{1}}\right) e^{G_{10} \gamma_{0} + G_{11} \gamma_{1}}\right)^{2}}
\end{equation*}

\noindent Residuals, each asserted to vanish:

\begin{center}\small
\begin{tabular}{lll}
\toprule
claim & residual & result\\
\midrule
d2/dbeta0 dbeta0 & $\partial^2_{\beta_0\beta_0} - \textrm{closed form}$ & $0$\\
d2/dbeta0 dbeta1 & $\partial^2_{\beta_0\beta_1} - \textrm{closed form}$ & $0$\\
d2/dbeta0 dgamma0 & $\partial^2_{\beta_0\gamma_0} - \textrm{closed form}$ & $0$\\
d2/dbeta0 dgamma1 & $\partial^2_{\beta_0\gamma_1} - \textrm{closed form}$ & $0$\\
d2/dbeta1 dbeta1 & $\partial^2_{\beta_1\beta_1} - \textrm{closed form}$ & $0$\\
d2/dbeta1 dgamma0 & $\partial^2_{\beta_1\gamma_0} - \textrm{closed form}$ & $0$\\
d2/dbeta1 dgamma1 & $\partial^2_{\beta_1\gamma_1} - \textrm{closed form}$ & $0$\\
d2/dgamma0 dgamma0 & $\partial^2_{\gamma_0\gamma_0} - \textrm{closed form}$ & $0$\\
d2/dgamma0 dgamma1 & $\partial^2_{\gamma_0\gamma_1} - \textrm{closed form}$ & $0$\\
d2/dgamma1 dgamma1 & $\partial^2_{\gamma_1\gamma_1} - \textrm{closed form}$ & $0$\\
\bottomrule
\end{tabular}
\end{center}

\noindent\textbf{Verdict:} every residual is exactly zero.

% Generated by cbmr-proofs. Do not edit.
% Every expression below is rendered from the object sympy verified, not transcribed.

\subsection{Covariance strategies}

Let $\mathrm{supp}(p)=\{c: L_{pc}\neq0\}$ and let $\sim$ be the transitive
closure of $c\sim c' \iff \exists p:\,c,c'\in\mathrm{supp}(p)$. Every distribution's Hessian has
$\beta\beta$ block $\sum_p L_{pc}L_{pc'}(B^\top\Sigma_pB)_{kk'}$, so the factor $L_{pc}L_{pc'}$
decides which entries can be nonzero at all.


\noindent Residuals, each asserted to vanish:

\begin{center}\small
\begin{tabular}{lll}
\toprule
claim & residual & result\\
\midrule
columns sharing no pattern have a zero cross block & $H_{(0k),(2k')}\ \textrm{with}\ 0\not\sim 2$ & $0$\\
characteristic polynomial factorises over the blocks & $\det(\mathrm{blkdiag}(A,B)-\lambda I)-\det(A-\lambda I)\det(B-\lambda I)$ & $0$\\
symmetric H satisfies H\^{}T H = H\^{}2 & $H^\top H - H^2$ & $0$\\
the singular values are the absolute eigenvalues & $\det(H^2-\lambda^2I)-\det(H-\lambda I)\det(H+\lambda I)$ & $0$\\
blockwise inverse & $\mathrm{blkdiag}(A,B)^{-1}-\mathrm{blkdiag}(A^{-1},B^{-1})$ & $0$\\
D.TL + C.BL = I & $D\,\mathrm{TL}+C\,\mathrm{BL}-I_m$ & $0$\\
D.TR + C.BR = 0 & $D\,\mathrm{TR}+C\,\mathrm{BR}$ & $0$\\
C\^{}T.TL + E.BL = 0 & $C^\top\mathrm{TL}+E\,\mathrm{BL}$ & $0$\\
C\^{}T.TR + E.BR = I & $C^\top\mathrm{TR}+E\,\mathrm{BR}-I_q$ & $0$\\
\bottomrule
\end{tabular}
\end{center}

\noindent\textbf{Verdict:} every residual is exactly zero.



\end{document}
